% TOP_PROBLEM: {"id": 4700026, "title": "A moments problem I", "category_id": 4, "category": {"id": 4, "name": "algebra", "display_name": "Algebra", "description": "Group theory, ring theory, field theory, and algebraic structures.", "slug": "algebra", "order_index": 4, "created_at": "2026-07-31T15:26:25.671Z"}, "status": "open", "problem_number": "AMR-046-0026", "set_id": 13}
% FIRST_POSTED: 2026-10-02
% ENTRY_KIND: statement_only
% UnsolvedMath ID 4700026; source catalogue code AMR-046-0026
% !TeX program = lualatex
% Definition and sources only; this file is not an LLM solution attempt.
\documentclass[11pt,a4paper]{article}
\usepackage[margin=25mm]{geometry}
\usepackage{fontspec}
\setmainfont{lmroman10-regular.otf}[BoldFont=lmroman10-bold.otf,
 ItalicFont=lmroman10-italic.otf,BoldItalicFont=lmroman10-bolditalic.otf]
\usepackage{amsmath,amssymb,mathtools}
\usepackage{unicode-math}
\setmathfont{latinmodern-math.otf}
\AtBeginDocument{\let\setminus\smallsetminus}
\usepackage{xurl,hyperref}
\hypersetup{colorlinks=true,urlcolor=blue}
\setcounter{secnumdepth}{0}
\setlength{\parindent}{0pt}
\setlength{\parskip}{5pt}
\setlength{\emergencystretch}{3em}
\begin{document}
\section*{A moments problem I}\label{4700026}
{\small UnsolvedMath ID 4700026 · AMR-046-0026 · Algebra · AMR Open Problem Lists}
\subsection{Definitions and mathematical statement}
Fix an integer $d\ge1$ and a polynomial
$f\in\mathbb C[x_1,\ldots,x_d]$. For every positive integer $m$, define
its complex moment over the real unit cube by
\[
 M_m(f)=\int_{[0,1]^d}f(x_1,\ldots,x_d)^m\,dx_1\cdots dx_d.
\]
Integration is with respect to ordinary product Lebesgue measure, applied
to real and imaginary parts. No complex conjugation or absolute value is
inserted in the integrand. The problem is whether
\[
 \bigl(M_m(f)=0\text{ for every }m\ge1\bigr)
       \quad\Longrightarrow\quad f\equiv0
\]
holds in every number of variables.

The conclusion means that every coefficient vanishes, equivalently that
the polynomial vanishes identically on the cube. The degree and coefficients
are unrestricted. Moments begin at one, because the zeroth moment equals
the volume of the cube and is always one. If $f$ were real-valued, its
second moment would immediately settle the question; the possible
cancellation between complex values is the essential difficulty.

The one-variable implication is known, as noted by the source, whereas
the question asks for arbitrary $d$. A uniform finite bound on the number
of moments needed is not assumed here. Nor does vanishing of integrals
against a restricted selection of test polynomials replace the required
vanishing of every positive integral power of the same polynomial. The
measure and the domain are fixed and are part of the formulation.
\subsection{Short English statement}
If every positive power of a complex polynomial integrates to zero over a real unit cube, must the polynomial vanish?
\subsection{Sources}
\begin{itemize}
\item[S1] ULAM.AI and the UnsolvedMath Contributors, supplied problems.json, record 4700026 (AMR-046-0026), AMR Open Problem Lists. Catalogue identity and source transcription; CC BY 4.0.
\url{https://huggingface.co/datasets/ulamai/UnsolvedMath}.
\item[S2] Gasull - Some open problems in low dimensional dynamical systems (2020), Problem environment 26: A moments problem I. Original source indexed by AMR-046-0026.
\url{https://arxiv.org/abs/2012.02524}.
\end{itemize}
\end{document}
